\documentclass[11pt]{article}
\usepackage[margin=1in]{geometry}
\usepackage{amsmath,amssymb,amsthm}
\usepackage{graphicx}
\usepackage{microtype}
\usepackage[hidelinks]{hyperref}

\newtheorem{theorem}{Theorem}
\newtheorem{lemma}{Lemma}
\newcommand{\cai}{\mathrm{c.a.i.}}
\newcommand{\otmax}{\mathbin{\otimes_{\max}}}

\title{Maximal Tensor Products Preserve Operator-Algebra Quotients\\
\large A candidate full answer to a question in arXiv:0804.0028}
\author{Automated Mathematics Run \texttt{fa\_banach\_001}, lane 4}
\date{27 August 2026}

\begin{document}
\maketitle

\begin{center}
\fbox{\parbox{0.9\textwidth}{\textbf{Status.}
Candidate full solution, likely valid; human review pending.}}
\end{center}

\section{Source question}

Blecher and Duncan study nuclearity-related properties of nonselfadjoint
operator algebras in \cite{BlecherDuncan}.  Immediately before their Lemma
2.7 (Lemma 2.8 in the arXiv source file), they write:

\begin{quote}
``To deal with quotients we will need the following.  We do not know if the
result is true for nonselfadjoint $D$.''
\end{quote}

Their lemma proves, when $D$ is a $C^*$-algebra, that an approximately unital
ideal $A\triangleleft B$ satisfies
\[
  (B\otmax D)/(A\otmax D)\cong (B/A)\otmax D
\]
completely isometrically.  It also proves the symmetric result when $B$ is a
$C^*$-algebra.  The question is whether selfadjointness of either tensor
factor is needed.  The published erratum \cite{Erratum} does not alter this
lemma or question.

\begin{figure}[ht]
  \centering
  \includegraphics[width=\textwidth]{figures/open_problem_crop.png}
  \caption{The question and Lemma 2.7 on page 4 of the source PDF.}
\end{figure}

\section{Definitions}

An operator algebra is a norm-closed, not necessarily selfadjoint algebra of
operators on a Hilbert space.  It is \emph{approximately unital} if it has a
contractive two-sided approximate identity.  All ideals below are closed and
two-sided.

The maximal operator-algebra tensor product $E\otmax F$ is characterized by
commuting completely contractive representations: if
$\pi:E\to B(H)$ and $\theta:F\to B(H)$ are commuting completely contractive
homomorphisms, then
\[
 e\otimes f\longmapsto \pi(e)\theta(f)
\]
extends completely contractively to $E\otmax F$.  Conversely, representations
of the maximal tensor product have this factor form, with the usual passage to
essential subspaces in the nonunital case; see \cite[Chapter 6]{BLM}.

An isomorphism of operator algebras is \emph{completely isometric} when it is
isometric on every matrix level.

\section{Proof intuition}

The quotient map clearly induces a complete contraction from the left-hand
quotient to $(B/A)\otmax D$.  The issue is the reverse matrix-norm inequality.
Represent the left-hand quotient completely isometrically.  Maximality splits
the lifted representation into commuting factors $\pi$ on $B$ and $\theta$
on $D$.  Since tensors from $A\otimes D$ vanish in the quotient,
$\pi(a)\theta(d)=0$.  On the essential subspace generated by $\theta(D)$ this
forces $\pi(a)=0$, so $\pi$ factors through $B/A$.  The two resulting factor
representations then give a representation of $(B/A)\otmax D$ which recovers
the original quotient representation without losing any matrix norm.

The essential-subspace step is the only point requiring care: it removes any
hidden nondegeneracy assumption and is precisely what allows $D$ to be
nonselfadjoint.

\section{The result}

\begin{theorem}\label{thm:main}
Let $A$ be an approximately unital ideal in an operator algebra $B$, and let
$D$ be an approximately unital operator algebra.  Then the canonical map
\[
 \Phi:\frac{B\otmax D}{A\otmax D}\longrightarrow (B/A)\otmax D,
 \qquad
 \Phi\bigl((b\otimes d)+(A\otmax D)\bigr)=(b+A)\otimes d,
\]
is a completely isometric isomorphism.  In particular,
\[
 0\longrightarrow A\otmax D\longrightarrow B\otmax D
 \longrightarrow (B/A)\otmax D\longrightarrow 0
\]
is an exact sequence of operator algebras.
\end{theorem}

\begin{proof}
First, $A\otmax D$ embeds completely isometrically in $B\otmax D$.
For completeness, this is the ideal-extension argument of
\cite[Proposition 2.5]{BlecherDuncan}.  If
$\pi:A\to B(H)$ and $\theta:D\to B(H)$ are nondegenerate commuting
completely contractive representations, the standard extension of $\pi$ from
the approximately unital ideal $A$ to $B$ is the completely contractive
homomorphism $\widetilde\pi$ determined on the dense subspace $\pi(A)H$ by
\[
 \widetilde\pi(b)\pi(a)\xi=\pi(ba)\xi.
\]
For $a\in A$, $b\in B$, $d\in D$, and $\xi\in H$,
\[
 \theta(d)\widetilde\pi(b)\pi(a)\xi
 =\theta(d)\pi(ba)\xi
 =\pi(ba)\theta(d)\xi
 =\widetilde\pi(b)\theta(d)\pi(a)\xi.
\]
Thus $\widetilde\pi$ commutes with $\theta$.  The universal property of
$\otmax$ gives the asserted complete isometry.  Its image is a closed ideal.

Put
\[
 X=B\otmax D,\qquad I=A\otmax D,\qquad E=X/I,
\]
and let $p:X\to E$ and $q:B\to B/A$ be the quotient maps.  Functoriality of
the maximal tensor product makes $q\otimes\mathrm{id}_D$ a completely
contractive homomorphism from $X$ to $(B/A)\otmax D$.  It annihilates $I$ and
therefore induces the completely contractive map $\Phi$ in the theorem.

It remains to prove that $\Phi$ is completely isometric.  Choose a completely
isometric representation $\rho:E\to B(H)$.  By the factor-representation
property of the maximal tensor product, there are commuting completely
contractive homomorphisms
\[
 \pi:B\to B(H),\qquad \theta:D\to B(H)
\]
such that
\begin{equation}\label{eq:factor}
 \rho\bigl(p(b\otimes d)\bigr)=\pi(b)\theta(d)
 \qquad (b\in B,\ d\in D).
\end{equation}
Let
\[
 K=\overline{\theta(D)H}.
\]
The space $K$ is invariant for $\pi(B)$ because $\pi(B)$ commutes with
$\theta(D)$.  Moreover, for $a\in A$ and $d\in D$, equation
\eqref{eq:factor} gives
\[
 \pi(a)\theta(d)=\rho\bigl(p(a\otimes d)\bigr)=0.
\]
Since $\theta(D)H$ is dense in $K$, it follows that
$\pi(a)|_K=0$ for every $a\in A$.  Consequently $\pi|_K$ factors through a
completely contractive homomorphism
\[
 \overline\pi:B/A\to B(K),\qquad
 \overline\pi(b+A)=\pi(b)|_K.
\]
The factor is completely contractive because $q:B\to B/A$ is a complete
quotient map.  Also $\theta_K(d)=\theta(d)|_K$ is a completely contractive
representation of $D$ on $K$, and it commutes with $\overline\pi(B/A)$.
Maximality therefore gives a completely contractive representation
\[
 \sigma:(B/A)\otmax D\to B(K),\qquad
 \sigma((b+A)\otimes d)=\overline\pi(b+A)\theta_K(d).
\]

For an algebraic tensor $x\in B\otimes D$, equations above show
\begin{equation}\label{eq:recover}
 \sigma\bigl(\Phi(p(x))\bigr)=\rho(p(x))|_K.
\end{equation}
Every summand $\pi(b)\theta(d)$ vanishes on $K^\perp$: if $(e_t)$ is a
contractive approximate identity for $D$, then
$\theta(e_t)$ converges strongly to the projection onto $K$, and
$\theta(d)=\theta(d)P_K$.  Its range is contained in $K$, since
$\pi(b)\theta(d)=\theta(d)\pi(b)$.  Hence restriction to $K$ does not change
the norm of $\rho(p(x))$, and the same assertion holds at every matrix level.
By density, \eqref{eq:recover} holds for every $u\in E$.  Therefore, for every
$n$ and $u\in M_n(E)$,
\[
 \|u\|
 =\|\rho^{(n)}(u)\|
 =\|\sigma^{(n)}(\Phi^{(n)}(u))\|
 \le \|\Phi^{(n)}(u)\|.
\]
The reverse inequality follows from complete contractivity of $\Phi$.
Thus $\Phi$ is completely isometric.

Finally, the range of $\Phi$ is closed and contains the algebraic tensor
product $(B/A)\otimes D$, because $q$ is onto.  That algebraic tensor product
is dense in $(B/A)\otmax D$, so $\Phi$ is onto.
\end{proof}

\section{Verification notes}

\begin{itemize}
\item The proof is matrix-level throughout: complete contractivity of
$\rho$, $\overline\pi$, $\theta_K$, and $\sigma$ gives the inequality in
$M_n$ for every $n$.
\item Degenerate representations cause no loss.  The tensor representation
is zero on $K^\perp$ and has range in $K$, so the compression used in the
proof preserves norms exactly.
\item No positivity or involution is used.  The only properties of $D$ are
the maximal-tensor universal property and the existence of a contractive
approximate identity.
\item There is no computational component to verify; the claim is purely
structural.
\end{itemize}

\section{Scope, novelty check, and review recommendation}

Theorem \ref{thm:main} gives the affirmative extension requested in the
source: both $B$ and $D$ may be nonselfadjoint.  Approximate unitality of $D$
is the standing class used by the source and is needed for the clean
essential-subspace argument.

A bounded search on 27 August 2026 covered the exact source title and authors,
the phrase ``nonselfadjoint $D$'', close variants of ``maximal tensor product
quotient/projective'', the source's citations, and the published erratum.  It
found Han's projectivity theorem for the maximal \emph{operator-system}
tensor product \cite{Han}, but no later paper explicitly answering this
operator-algebra question and no exact occurrence of Theorem \ref{thm:main}.
Novelty is plausible rather than certified.

Human review should focus on the standard factor-representation property in
the precise convention for $\otmax$ used in \cite{BlecherDuncan}.  Once that
defining property is confirmed, the remainder is an elementary quotient and
essential-subspace argument.

\begin{thebibliography}{9}

\bibitem{BlecherDuncan}
D. P. Blecher and B. L. Duncan,
\emph{Nuclearity-related properties for nonselfadjoint algebras},
J. Operator Theory \textbf{65} (2011), 47--70;
arXiv:0804.0028.

\bibitem{Erratum}
D. P. Blecher and B. L. Duncan,
\emph{Errata to ``Nuclearity-related properties for nonselfadjoint
algebras''}, J. Operator Theory \textbf{65} (2011), 471--473.

\bibitem{BLM}
D. P. Blecher and C. Le Merdy,
\emph{Operator Algebras and Their Modules: An Operator Space Approach},
London Mathematical Society Monographs 30, Oxford University Press, 2004.

\bibitem{Han}
K. H. Han,
\emph{On maximal tensor products and quotient maps of operator systems},
J. Math. Anal. Appl. \textbf{384} (2011), 375--386;
arXiv:1010.0380.

\end{thebibliography}

\end{document}
